%% =========================================================================
%% 01-fundamentals.tex -- Section 1.
%%
%% Scope: the section heading, its opening paragraph, 1.1 Spiking Neuron
%% Model and 1.2 Spiking Convolution (three-dimensional, with Algorithm 1).
%%
%% The text states the model only. It carries no implementation numbers
%% (precision, widths, clock rates) and no hardware argument (throughput,
%% cost, accuracy trade-offs); those belong to the sections that describe
%% the IP.
%%
%% Notation: channel counts are \Cin / \Cout, channel indices \cin / \cout.
%% Depth is D with subscripts i/o; the window is described by kernel size K,
%% stride S and padding P, which is the configuration the report targets.
%%
%% Build through main.tex (IEEEtran, two columns):
%%   latexmk -pdf -interaction=nonstopmode main.tex
%% =========================================================================

\section{SNN Computation Fundamentals}
\label{sec:1-fund}

A spiking neural network (SNN) carries information in binary spikes emitted at
discrete time steps, rather than in real-valued activations propagated in a
single pass. Each neuron integrates the weighted spikes that reach it into a
membrane potential, lets that potential decay according to the leakage of the
cell, and emits a spike of its own whenever the potential crosses a threshold,
after which the potential is reset. A spiking convolutional neural network
(SCNN), the form used in this report, arranges such neurons as a convolutional
layer: the fan-in of a neuron is then a sliding window over an input spike
volume, and the layer repeats a weighted sum over time steps, input channels and
kernel positions. The rest of this section states these definitions precisely
--- the neuron dynamics, the variables of a layer, and the order in which they
are evaluated.

\subsection{Spiking Neuron Model}
\label{sec:1-neuron}

The neuron is the only stateful element of the network, which means that
everything the network must remember between two time steps is the set of
membrane potentials; currents and spikes are derived at each step from that
state and from the input. At step $t$ it receives a current from its presynaptic
neighbours, updates its membrane potential, and decides whether to emit a spike;
the spike is binary and is the only signal the neuron sends. The current of that
step is the weighted sum of the spikes arriving at it, plus a bias,
\begin{equation}
I_i[t] = \sum_{j} w_{ij}\,s_j[t] + b_i ,
\label{eq:1-current}
\end{equation}
where $s_j[t]\in\{0,1\}$ is the spike of presynaptic neuron $j$ and $w_{ij}$ the
synapse that connects it to neuron $i$.

In continuous form the membrane obeys
\begin{equation}
\tau_m \frac{dv_i}{dt} = -v_i + R\,I_i(t),
\label{eq:1-ode}
\end{equation}
where $\tau_m$ is the membrane time constant, which sets how fast the potential
decays, and $R$ is the membrane resistance, which converts the input current
into a change of potential. A
digital simulation evaluates this dynamics once per step and carries the
potential of the previous step forward with a decay factor,
\begin{equation}
v_i[t] = \left(1-\frac{1}{\tau_m}\right)u_i[t] + I_i[t] .
\label{eq:1-membrane}
\end{equation}
Equation~\eqref{eq:1-membrane} is an accumulator with forgetting: the potential
left by the previous step is kept, reduced by a fraction $1/\tau_m$, and the
current of the present step is added. A small $\tau_m$ makes the neuron forget
quickly, so only recent input counts; a large $\tau_m$ lets earlier input weigh
for longer. When $\tau_m$ is infinite the decay disappears and the neuron simply
sums its input over time, which is the integrate-and-fire (IF) model.

The output of the step is a comparison, and the potential is rewritten at the
same time:
\begin{equation}
s_i[t+1] =
\begin{cases}
1, & v_i[t] \geq V_{th}, \\
0, & v_i[t] < V_{th},
\end{cases}
\label{eq:1-spike}
\end{equation}
\begin{equation}
u_i[t+1] = v_i[t] - s_i[t+1]\,V_{th} ,
\label{eq:1-reset}
\end{equation}
where $V_{th}$ is the firing threshold. What the neuron sends is this comparison
and nothing else: its activity is described by when and how often it fires, not
by the height of the potential. 

Table~\ref{tab:1-neurons} sets the two models side by side. They share the reset
rule of \eqref{eq:1-reset} and differ in one place only, the decay term of
\eqref{eq:1-membrane}: the IF neuron integrates without loss, the leaky
integrate-and-fire (LIF) neuron lets the potential relax between inputs.

\begin{table}[!ht]
\centering
\caption{Spiking neuron models. Both subtract the threshold after a spike and
differ only in the decay of the membrane.}
\label{tab:1-neurons}
\footnotesize
\begin{tabularx}{\columnwidth}{|p{1.0cm}|X|p{2.3cm}|}
\hline
\textbf{Model} & \textbf{Membrane update} & \textbf{After a spike} \\
\hline
IF & $v_i[t]=u_i[t]+I_i[t]$ & $v_i[t]-V_{th}$ \\
\hline
LIF & $v_i[t]=\left(1-\frac{1}{\tau_m}\right)u_i[t]+I_i[t]$ & $v_i[t]-V_{th}$ \\
\hline
\end{tabularx}
\end{table}

A variant resets the potential to rest instead, $u_i[t+1]=0$. Subtracting
$V_{th}$ leaves a higher potential for the next step, so the neuron sits closer
to the threshold and fires more readily than after a reset to rest.

Two properties of this model are used in what follows. The state of a layer at
step $t$ is fully described by the membrane potentials of its neurons, because
the spikes of that step are a function of that state and of the input. And the
output of a layer is defined only once its $T$ steps have been evaluated in
order, since each step consumes the state left by the previous one.

\subsection{Spiking Convolution}
\label{sec:1-conv}

\figplaceholderwide{../../figures/3DCNN-Architecture.pdf}{A three-dimensional
SCNN pipeline over $T=8$ time steps, drawn left to right on the lower row and
right to left on the upper one. A layer holds one spike volume of size
$C\times D\times H\times W$ per time step; the input here is
$1\times3\times24\times24$. Each $3$-D convolutional layer applies a
$7\times7\times7$ kernel at unit stride with same padding, which preserves $D$,
$H$ and $W$ and changes only the channel count, from $1$ to $32$, $64$ and
$128$; a $2\times2$ pooling then halves the map, and global average pooling with
a fully connected layer yields $8$ outputs per time step. The weights are shared
by all $T$ steps, so the parameter count does not grow with $T$.}{fig:1-pipeline}

A layer of a spiking convolutional neural network places its neurons on a
three-dimensional lattice, one channel of neurons per lattice position, and
connects each neuron to a small neighbourhood of the previous layer rather than
to all of it. The spikes produced by such a layer are a function of the time
step, the input channel and the position in the volume,
$s(t,c,d,h,w)\in\{0,1\}$, and this report describes them as $T$ binary volumes
of size $C\times D\times H\times W$, one per time step. A volume is the
description a convolution asks for, since a convolution is a statement about
where a spike sits relative to its neighbours.

With the volumes in hand, the input current of the neuron at output channel
$\cout$ and position $(d,h,w)$ is the weighted sum of the spikes inside its
window,
\begin{equation}
\begin{split}
i_{\cout}(d,h,w,t) &= \sum_{\cin}\sum_{\mathbf{k}\in\mathcal{K}}
w_{\cout,\cin,\mathbf{k}}\\
&\quad\times s(t,\cin,\; dS_d-P_d+k_d,\\
&\hspace{2.8em} hS_h-P_h+k_h,\; wS_w-P_w+k_w),
\end{split}
\label{eq:1-conv}
\end{equation}
where $\mathcal{K}=\{0,\dots,K_d-1\}\times\{0,\dots,K_h-1\}\times\{0,\dots,K_w-1\}$
is the set of kernel positions and $\mathbf{k}=(k_d,k_h,k_w)$ is one of them;
$S_d,S_h,S_w$ are the strides and $P_d,P_h,P_w$ the paddings; and a position
that falls outside the input volume contributes nothing.
Equation~\eqref{eq:1-conv} is the same three-dimensional convolution as in a
non-spiking network; what singles out the spiking case is that the operand $s$
is binary, so every term is either the weight itself or zero, and the sum is a
weighted count of the spikes that fall inside the window.

The window visits every output position of the volume, every input channel and
every time step, which fixes the extent of the output volume along each
geometric axis,
\begin{equation}
D_o = \left\lfloor \frac{D_i + 2P_d - K_d}{S_d} \right\rfloor + 1 ,
\label{eq:1-dout}
\end{equation}
and likewise for $H_o$ and $W_o$ with the parameters of the other two axes.
Figure~\ref{fig:1-pipeline} shows one such layer inside a network: every layer
holds $T$ volumes of $C\times D\times H\times W$ spikes, the window of
\eqref{eq:1-conv} spans $K_d\times K_h\times K_w$ positions inside each of them,
and the same weights are applied at every time step.

Two properties follow from the weight sharing. The number of parameters of a
layer is $\Cout \Cin K_d K_h K_w$ --- one weight for each combination of output
channel, input channel and kernel position --- independent of the size of the
volume and of the number of time steps, so a layer costs the same to store
whether it runs for $T=4$ or $T=64$ steps. And a spike affects only the
$K_d \times K_h \times K_w$ positions of the output volume that contain it in
their window, and only at the time step at which it occurs, so the layer is
local in both space and time.

In total, one pass over a layer performs
\begin{equation}
N_{\mathrm{syn}} = T \cdot \Cout D_o H_o W_o \cdot \Cin K_d K_h K_w
\label{eq:1-count}
\end{equation}
synaptic operations, each of which adds a weight or adds nothing. The count
grows linearly with the number of time steps, whereas the parameter count does
not.

Algorithm~\ref{alg:1-layer} writes \eqref{eq:1-conv} out as a loop nest. The
volumes $S$ have size $T\times\Cin\times D_i\times H_i\times W_i$, the weights
are indexed by output channel, input channel and kernel position, and $U$ holds
the membrane potential of every output neuron; no step overlaps with another.

The order of the loops is a choice. Any permutation that keeps the temporal loop
outside the others returns the same output, because the only state carried from
one step to the next is the membrane potential, and the spikes of step $t$ are
produced after the whole fan-in of step $t$ has been summed. The order above is
kept as the reference: it states what a layer computes with the fewest
assumptions, and every arrangement discussed later is a reordering of these same
loop bodies.

\begin{algorithm}[!htb]
\caption{One SCNN layer, evaluated sequentially.}
\label{alg:1-layer}
\scriptsize
\algrenewcommand\algorithmicindent{1.1em}
\begin{algorithmic}[1]
\Require $S(T,\Cin,D_i,H_i,W_i)$, $W(\Cout,\Cin,K_d,K_h,K_w)$
\Statex \hspace{2.4em}$b(\Cout)$, $V_{th}$, $U(\Cout,D_o,H_o,W_o)$
\Ensure $O(T,\Cout,D_o,H_o,W_o)$ and the updated $U$
\For{$t \gets 0$ \textbf{to} $T-1$}
  \For{each $(\cout,d,h,w)$ of the output volume}
    \State $i \gets b[\cout]$
    \For{$\cin \gets 1$ \textbf{to} $\Cin$}
      \For{$(k_d,k_h,k_w)\in\mathcal{K}$}
        \State $d' \gets dS_d-P_d+k_d$
        \State $h' \gets hS_h-P_h+k_h$;\quad $w' \gets wS_w-P_w+k_w$
        \If{$(d',h',w')$ lies inside the input volume}
          \State $i \gets i + W[\cout,\cin,k_d,k_h,k_w]\,S[t,\cin,d',h',w']$
        \EndIf
      \EndFor
    \EndFor
    \State $v \gets (1-1/\tau_m)\,U[\cout,d,h,w] + i$
    \If{$v \ge V_{th}$}
      \State $O[t,\cout,d,h,w] \gets 1$;\quad $U[\cout,d,h,w] \gets v-V_{th}$
    \Else
      \State $O[t,\cout,d,h,w] \gets 0$;\quad $U[\cout,d,h,w] \gets v$
    \EndIf
  \EndFor
\EndFor
\end{algorithmic}
\end{algorithm}

\subsection{What the Sequential Model Costs}
\label{sec:1-cost}

Algorithm~\ref{alg:1-layer} is a definition of what a layer computes, and it is
also a schedule for computing it, and the schedule has a price. Three quantities
describe that price, and they are not interchangeable: the arithmetic a network
performs, the state it has to keep while performing it, and the length of the
chain of dependencies into which the arithmetic is arranged.

The arithmetic is the sum of \eqref{eq:1-count} over the layers,
\begin{equation}
W = \sum_{\ell=1}^{N_{\mathrm{lay}}} N_{\mathrm{syn}}^{(\ell)} ,
\label{eq:1-work}
\end{equation}
where the superscript indexes the layer. This is a bound and not a count,
because a term of \eqref{eq:1-conv} contributes the weight only when the spike is
one: what a step actually adds is proportional to the number of spikes that enter
it, and each of those spikes costs one weight addition for every output channel
and kernel position it can reach,
\begin{equation}
\tilde{W}(t) = \Cout K_d K_h K_w \, A(t) ,
\label{eq:1-workt}
\end{equation}
where $A(t)$ counts the spikes arriving at step $t$. If a fraction $\rho$ of the
input spikes are ones, then $\mathbb{E}[\tilde{W}] = \rho W$: the work of a layer
is a property of its input, and \eqref{eq:1-work} is the value it takes when every
spike is one.

The state is the set of membrane potentials that must survive from one step to
the next. By the first property of Section~\ref{sec:1-neuron} that set is the
whole of what a layer remembers, and its size is
\begin{equation}
M = \max_{\ell}\; \Cout^{(\ell)} D_o^{(\ell)} H_o^{(\ell)} W_o^{(\ell)} ,
\label{eq:1-state}
\end{equation}
the neuron count of the largest layer. It does not grow with $T$: the same
potentials are carried through every step, and each step overwrites them.

The chain is the quantity that does not shrink. A step has to accumulate its
fan-in before it can compare the result with $V_{th}$, and the reset of
\eqref{eq:1-reset} makes the next step depend on that comparison, so the steps of
a layer form a chain of length $T$ and the chain of the network is
\begin{equation}
P = T \sum_{\ell=1}^{N_{\mathrm{lay}}} d_\ell ,
\label{eq:1-path}
\end{equation}
where $d_\ell$ is the depth of one step of layer $\ell$: the fan-in accumulation
of \eqref{eq:1-conv} followed by the comparison of \eqref{eq:1-spike}. The factor
$T$ multiplies rather than adds because the steps are chained rather than merely
repeated; step $t+1$ cannot update its membrane before step $t$ has produced its
spikes.

The three quantities depend on $T$ differently, and the difference carries the
rest of the argument. $M$ does not depend on it at all, while $W$ and $P$ both
grow linearly with it, $W$ because $T$ more steps are executed and $P$ because
$T$ more steps are chained. Sparsity separates them further: it reduces
$\tilde{W}$ and leaves $P$ untouched, because a step in which no spike arrives is
still a step. The potential of \eqref{eq:1-membrane} decays by $1/\tau_m$ in it
and the comparison of \eqref{eq:1-spike} is still made, so an empty step changes
the state and cannot be dropped without changing the answer.

Nothing in the definitions so far asks whether all $T$ steps are needed. The
output is defined by \eqref{eq:1-count} over the whole interval, so the cost of a
layer is fixed when $T$ is chosen and not when the input arrives. Yet the second
property of Section~\ref{sec:1-neuron} says that the state at step $t$ is a
complete description of everything the network has computed up to that step, and
the question whether the output after $t$ steps is the intended one is therefore
well posed at every $t$. The sequential model does not ask it; it evaluates the
remaining steps as if the answer could not change. Whether a partial output is
sufficient is a property of the trained network and not of these definitions,
which is why the question has to be settled outside them.

What the sequential model costs is then a chain of $T$ stages that no reordering
of the loops removes, and the state of the largest layer, which no step reduces.
What it does not do is ask, at any point along that chain, whether the answer is
already the one that would be reported.
